% Part: conditional-logics
% Chapter: minimal-change-semantics
% Section: true-false

\documentclass[../../../include/open-logic-section]{subfiles}

\begin{document}

\olfileid[zh]{con}{min}{tf}

\olsection{反事实条件句的真与假}

若最接近的 $!A$ 世界都是 $!B$ 世界，则反事实条件句 $!A \cif !B$（非空真地）为真，如 \olref{fig:true} 所示。
\begin{figure}
\begin{center}
\begin{tikzpicture}[scale=.7]\small
  \spheresystem{5}
  \draw (0,0) node {$w$};
  \propositionintersect{0}{4}{40}{3.5}
  {\tikzset{proposition/.append={smooth,tension=1.4}}
  \proposition[shift={(1.6,-1)}]{90}{1}{30}{4}}
  \path (1.5,3) node[below] {$\formula{B}$};
  \path (3,0) node {$\formula{A}$};
\end{tikzpicture}
\caption{非空真地为真的反事实条件句}
\ollabel{fig:true}
\end{center}
\end{figure}
如果围绕 $w$ 的球面系统根本没有允许 $!A$ 的球面，那么反事实条件句在 $w$ 处也为真。在此情形下，它是\emph{空真地}为真的（参见 \olref{fig:vacuous}）。
\begin{figure}
\begin{center}
\begin{tikzpicture}[scale=.7]\small
  \spheresystem{5}
  \draw (0,0) node {$w$};
  \proposition{0}{6.5}{40}{3.5}
  {\tikzset{proposition/.append={smooth,tension=1.4}}
  \proposition[shift={(1.2,-1)}]{90}{1}{30}{4}}
  \path (1.5,3) node[below] {$\formula{B}$};
  \spherepos{0}{7}{node {$\formula{A}$}}
\end{tikzpicture}
\caption{空真地为真的反事实条件句}
\ollabel{fig:vacuous}
\end{center}
\end{figure}

它可以以两种方式为假。一种方式是，最接近的 $!A$ 世界并非都是 $!B$ 世界，但其中有些是 $!B$ 世界。在此情形下，$!A \cif
\lnot !B$ 也为假（参见 \olref{fig:false}）。
\begin{figure}
\begin{center}
\begin{tikzpicture}[scale=.7]\small
  \spheresystem{5}
  \draw (0,0) node {$w$};
  \propositionintersect{0}{4}{40}{3.5}
  {\tikzset{proposition/.append={smooth,tension=1.4}}
  \proposition[shift={(1.6,0)}]{90}{1}{30}{3}}
  \path (1.5,3) node[below] {$\formula{B}$};
  \path (3,0) node {$\formula{A}$};
\end{tikzpicture}
\caption{反事实条件句为假，其颠倒式亦为假}
\ollabel{fig:false}
\end{center}
\end{figure}
如果最接近的 $!A$ 世界与 $!B$ 世界完全不重叠，那么 $!A \cif !B$ 为假。但在这种情形下，所有最接近的 $!A$ 世界都是
$\lnot !B$ 世界，因此 $!A \cif \lnot !B$ 为真（参见
\olref{fig:false-opposite}）。
\begin{figure}
\begin{center}
\begin{tikzpicture}[scale=.7]\small
  \spheresystem{5}
  \draw (0,0) node {$w$};
  \propositionintersect{0}{4}{40}{3.5}
  {\tikzset{proposition/.append={smooth,tension=1.4}}
  \proposition[shift={(-1,0)}]{90}{1}{30}{3}}
  \path (-1,3) node[below] {$\formula{B}$};
  \path (3,0) node {$\formula{A}$};
\end{tikzpicture}
\caption{反事实条件句为假，其颠倒式为真}
\ollabel{fig:false-opposite}
\end{center}
\end{figure}

与严格条件句相比，反事实条件句可能是偶然的。考虑 \olref{fig:contingent} 中的球面模型。最接近 $u$ 的 $!A$ 世界都是 $!B$ 世界，因此 $\mSat{M}{!A \cif
  !B}[u]$。但最接近 $v$ 的 $!A$ 世界中有一些不是
$!B$ 世界，因此 $\mSat/{M}{!A \cif !B}[v]$。
\begin{figure}
\begin{center}
\begin{tikzpicture}[scale=.6]\small
  \spheresystem{7}
  \spheresystem[shift={(0:3.1)},dashed]{4}
  \propositionintersect{45}{4}{40}{4.5}
  \begin{scope}
    \clip \propositionplot{45}{4}{40}{4.5} ;
    \spherelayer[shift={(0:3.1)}]{3}
  \end{scope}
  \proposition[shift={(1.4,-.5)}]{90}{1}{35}{5}
  \path (0,0) node {$u$};
  \path (0:3.1) node {$v$};
  \spherepos{35}{9}{node {$\formula{A}$}}
  \spherepos{80}{9}{node {$\formula{B}$}}
\end{tikzpicture}
\end{center}
\caption{偶然的反事实条件句}
\ollabel{fig:contingent}
\end{figure}
\end{document}
